% ===================================================================
%  TABELO N-RO 11 — La urbo kaj ĝiaj monumentoj. — La incendio.
%  Traduction 2026 d'apres texto/io, controlee sur texto/fr et
%  texto/es. Consigne : voir texto/eo/10-tabelo-01.tex.
% ===================================================================

%%K t11-apar-1 apar
\VUsaut{2.0mm}
\VUcentre{13.2pt}{90}{TRIA SERIO}
\VUsaut{3.0mm}
\VUfilet{16mm}
\VUsaut{5.6mm}
\VUtitre{15.0pt}{60}{TABELO N\textsuperscript{o} 11}
\VUsaut{1.4mm}
\VUpk{11.4pt}{40}{La urbo kaj ĝiaj monumentoj. --- La incendio.}
\VUsaut{2.2mm}

%%K t11-tit-1 sub
\VUpk{10.2pt}{40}{La antaŭurbo.}
\VUsaut{3.0mm}

%%K t11-c1-01-1 p
1. --- Mi loĝas grandan urbon situantan je 100 kilometroj de Oceano, kaj kiu
tamen estas \VUgras{haveno} \textsuperscript{(1)} altranga. La fluo kiu
borderas ĝin havas larĝon de 400 metroj kaj formas tre belan radon.

%%K t11-c1-02-1 p
2. --- La tuta parto de la urbo kiu estas sur la \VUgras{kajoj}
\textsuperscript{(2)} tute aspektas mara kaj industria, kaj formas
\VUgras{antaŭurbon} \textsuperscript{(3)} kiu estas loĝata nur de ŝipanoj kaj
laboristoj. Ĉi tiuj laboras en la \VUgras{fabrikejoj} \textsuperscript{(4)} kaj
en la \VUgras{gasfarejo} \textsuperscript{(5)} kies alta \VUgras{kameno}
\textsuperscript{(6)} kaj \VUgras{gasometro} videblas de tre malproksime.

%%K t11-c1-03-1 p
3. --- En la mezepoko, la urbo estis fortikigita, kaj oni ankoraŭ trovas
kelkajn restaĵojn de ĝiaj remparoj, aparte la \VUgras{pordon}
\textsuperscript{(8)} de la malnova \VUgras{fortikigita kastelo}
\textsuperscript{(7)} flankumitan de ĝiaj du \VUgras{turoj}
\textsuperscript{(9)} kiuj nun uzatas kiel \VUgras{karcero}.

%%K t11-c1-04-1 p
4. --- En ĉi tiu tuta \VUgras{kvartalo}, kiu aspektas tre malnova, nur unu
konstruaĵo estas relative moderna : tio estas la \VUgras{doganejo}
\textsuperscript{(10)}. Ĝi estas situanta inter la kajo kaj la \VUgras{strateto}
\textsuperscript{(11)} kiu kondukas al la malnova \VUgras{komunuma domo}
\textsuperscript{(12)}. Oni facile rekonas tiun lastan monumenton pro la
\VUgras{sonorilejo} \textsuperscript{(13)} giganta kiu superregas la antikvan
urbon.

%%K t11-c1-05-1 p
5. --- Sur la alia flanko de la strato, \VUgras{staras} edifico konstruita
samstile kiel la doganejo : tio estas la \VUgras{Borso} \textsuperscript{(14)}.
Unu el ĝiaj fasadoj rigardas malgrandan placon kie estas la
\VUgras{prefektejo} \textsuperscript{(15)}. Fakte, nia urbo sen esti ĉefurbego,
estas tamen grava ĉefurbo de departemento.

%%K t11-tit-2 sub
\VUsaut{5.0mm}
\VUpk{10.2pt}{40}{La plej alta parto de la urbo.}
\VUsaut{3.4mm}

%%K t11-c1-06-1 p
6. --- La garnizono, sufiĉe grandnombra, estas loĝata en vastaj
\VUgras{kazernoj} \textsuperscript{(16)} tre salubraj; ili etendas sin ĝis la
kradaro de la \VUgras{parko komunuma} \textsuperscript{(17)}. La
\VUgras{bulvardo} \textsuperscript{(19)} kiu kondukas al tiu parko pasas
malantaŭ la \VUgras{ŝparkasejo} \textsuperscript{(18)} kaj finiĝas ĉe la placo
de la \VUgras{katedralo} \textsuperscript{(20)}.

%%K t11-c1-07-1 p
7. --- Tiuj ĉiuj monumentoj sin etendas amfiteatre sur monteto kie estas ankaŭ
nia vasta \VUgras{liceo} \textsuperscript{(21)}.

%%K t11-c1-07-2 p
Se oni malsupreniras tra vojo iom kliniĝinta, kiu rekuniĝas kun la Granda
strato, oni alvenas al la \VUgras{Tribunalo} \textsuperscript{(22)} antaŭ kiu
etendas sin agrabla \VUgras{ĝardeno publika} \textsuperscript{(26)}. Tiu
\VUgras{ĝardenplaco} ne estas tre granda, sed ĝi faras meze de la urbo oazon de
verdaĵo kaj malvarmeto. Pro tio ĝi estas tre vizitata de la urbanoj kiuj ŝatas
veni sidiĝi sub la tendo de la \VUgras{kafejo} \textsuperscript{(25)}, por
aŭskulti la soldatajn muzikistojn kiuj ĵaŭde kaj dimanĉe ludas en la
\VUgras{kiosko} \textsuperscript{(27)} metita inter la herbejoj kaj masivoj.

%%K t11-c1-08-1 p
8. --- Sur unu el tiuj herbejoj staras bronza \VUgras{statuo}
\textsuperscript{(28)}, monumento starigita por la memoro de unu el niaj
samurbanoj, kiu faris grandajn servojn al sia naskiĝurbo.

%%K t11-tit-3 sub
\VUsaut{4.0mm}
\VUpk{10.2pt}{40}{La transportiloj.}
\VUsaut{3.4mm}

%%K t11-c1-09-1 p
9. --- Ĝis nun nia urbo ne jam estas lumigata per la lumo elektra, ĝi havas nur
\VUgras{gasbrulilojn} \textsuperscript{(29)}. Sed la \VUgras{ĉi-tiea ĵurnalaro}
kampanjas por ke tiu sistemo de lumigo estu plibonigata, \VUgras{kiel} oni jam
perfektigis \VUgras{la sistemon de trakcio de la} tramvojoj. \VUgras{La}
malnovaj veturiloj, \VUgras{tirataj} de \VUgras{ĉevaloj}, estis praktike
anstataŭigitaj per \VUgras{elektraj tramvojoj} \textsuperscript{(31)} kiuj
rapide transportas la vojaĝantojn de unu ekstremo de la \VUgras{urbo al la
alia}, po tre \VUgras{malkara} prezo.

%%K t11-c1-10-1 p
10. --- Apud la Juĝejo estas la \VUgras{Banko} \textsuperscript{(32)} kie oni
diskontas la komercajn valorojn. La \VUgras{bankkomisiistoj}
\textsuperscript{(33)} iras enkasigi hejme la ŝuldaĵojn.

%%K t11-tit-4 sub
\VUsaut{4.0mm}
\VUpk{10.2pt}{40}{Arto kaj instruo.}
\VUsaut{3.4mm}

%%K t11-c1-11-1 p
11. --- La bela edifico kiu staras proksime, estas la \VUgras{muzeo}. Ĝi kune
enhavas la \VUgras{bibliotekon} \textsuperscript{(34)} de la urbo, la belajn
\VUgras{kolektaĵojn pri naturscienco} \textsuperscript{(35)} (Naturmuzeo) kaj
gracian \VUgras{pentromuzeon} \textsuperscript{(36)} kiu konsistigas belegan
\VUgras{ekspoziciejon} permanentan.

%%K t11-c1-11-2 p
Ĝi malfermiĝas por la publiko dufoje ĉiusemajne, sed la fremduloj rajtas
viziti ĝin iun ajn tagon po gratifiketo al la \VUgras{gardisto}
\textsuperscript{(37)}. La \VUgras{pentristoj} \textsuperscript{(38)} venas tien
por labori. Oni vidas ilin \VUgras{antaŭ} sia trestlo, kun \VUgras{paletro} en
la mano, kopiante la famajn pentraĵojn.

%%K t11-c1-12-1 p
12. --- Sur la altaĵo, malantaŭ la muzeo, oni ekvidas la \VUgras{kupolon}
\textsuperscript{(40)} de la \VUgras{hospitalo} \textsuperscript{(39)} kaj la
konstruaĵojn de la \VUgras{normala lernejo} \textsuperscript{(41)} en kiu estis
ekzercataj la instruistoj de niaj lernejoj komunumaj. Sur la placo de la Teatro
estas \VUgras{poŝtoficejo} \textsuperscript{(43)} instalita en \VUgras{domo
privata} \textsuperscript{(42)}.

%%K t11-tit-5 sub
\VUsaut{5.0mm}
\VUpk{11.4pt}{40}{La incendio.}
\VUsaut{3.0mm}

%%K t11-tit-6 sub
\VUpk{10.2pt}{40}{Krio pri incendio.}
\VUsaut{3.4mm}

%%K t11-c2-01-1 p
1. --- Teruriga rumoro disvastiĝas tra la urbo : incendio ĵus okazis en la
teatro! En la momento kiam mi alvenas sur la \VUgras{placon}
\textsuperscript{(96)}, jam la homamaso estas kunigita sur la
\VUgras{trotuaro} \textsuperscript{(46)} kaj sur la \VUgras{ŝoseo}
\textsuperscript{(47)}. La \VUgras{poŝtistoj} \textsuperscript{(44)} kaj la
\VUgras{leteristoj} \textsuperscript{(45)} eliras el la poŝtoficejo.
\VUgras{Tramkondukisto} \textsuperscript{(48)} devis haltigi sian (tram)veturilon
por lasi pasi \VUgras{ambulancveturilon} \textsuperscript{(50)} urban, kiu
alvenas kun \VUgras{kirurgo} \textsuperscript{(52)} kaj \VUgras{flegisto}
\textsuperscript{(51)} por flegi \VUgras{vunditon} \textsuperscript{(53)}
alportitan per \VUgras{brankardo} \textsuperscript{(54)} apud la
\VUgras{ĵurnalkiosko} \textsuperscript{(55)}.

%%K t11-c2-02-1 p
2. --- Infanteria kompanio, kiu estis revenanta de la ekzercoj, haltis por
sekurigi la ordan servadon kun la \VUgras{rajdantaj ĝendarmoj urbaj}
\textsuperscript{(61)}. \VUgras{La rajdantoj}, kies ĉevaloj baŭmas, pli bone ol
la \VUgras{soldatoj} \textsuperscript{(57)} aŭ la \VUgras{ĝendarmoj}
\textsuperscript{(63)}, reirigas la \VUgras{scivolulojn}
\textsuperscript{(56)}. Cetere la ĝendarmoj estas tre okupataj. Unu el ili
montras al la \VUgras{policistoj} \textsuperscript{(64)} kien oni devas
transporti alian \VUgras{akcidentiton} \textsuperscript{(65)}, fajrobrigadanon
kuraĝan falintan de ŝtupetaro.

%%K t11-c2-03-1 p
3. --- \VUgras{La oficiro} \textsuperscript{(66)} precizigas la lokon kie oni
devas loki la \VUgras{vaporpumpilon} \textsuperscript{(67)} alvenantan.
Atendante tiun potencan maŝinon, li baterie instaligis \VUgras{manpumpilon}
\textsuperscript{(68)} kies longa \VUgras{tubo} \textsuperscript{(69)} lanĉas
akvon torente sur la fajron. Ses fajrobrigadanoj manovras ĝin, dum ilia
kamarado, kiu tenas la \VUgras{ŝprucilon} \textsuperscript{(70)}, direktas la
\VUgras{ŝprucon} \textsuperscript{(71)} al la centro de la incendio.

%%K t11-c2-04-1 p
4. --- Ĉe la alia flanko de la teatro, oni apogis la grandan
\VUgras{ŝtupetaron} \textsuperscript{(72)}. Ĝi ebligas al la fajrobrigadanoj
alproksimiĝi al la flamoj kiuj ŝprucas inter kirloj de nigra \VUgras{fumo}
\textsuperscript{(75)}.

%%K t11-tit-7 sub
\VUsaut{5.0mm}
\VUpk{10.2pt}{40}{Bravaj faroj.}
\VUsaut{3.4mm}

%%K t11-c2-05-1 p
5. --- \VUgras{La incendio} \textsuperscript{(74)} naskiĝis en la vestiblo de la
\VUgras{teatro} \textsuperscript{(73)}, dum oni ripetis \VUgras{la operon} kies
unua \VUgras{prezentado} estus okazinta morgaŭ. Feliĉe nur malmultaj
\VUgras{spektantoj} \textsuperscript{(76)} estis en la spektosalono. La
kompatinduloj rifuĝis sur la \VUgras{balkonon} \textsuperscript{(77)} kie
kuraĝaj \VUgras{fajrobrigadanoj} \textsuperscript{(86)} venas helpi ilin per
ŝtupetaro kaj ŝnuregoj.

%%K t11-c2-06-1 p
6. --- Sub la \VUgras{peristilo} \textsuperscript{(78)}, kie la \VUgras{afiŝo}
\textsuperscript{(79)} informas pri la prezentado, oni vidas la
\VUgras{muzikistojn} \textsuperscript{(81)} de \VUgras{la orkestro}
\textsuperscript{(80)} forkuri, kunportante siajn instrumentojn :
\VUgras{violono} \textsuperscript{(81)}, \VUgras{fluto} \textsuperscript{(82)},
\VUgras{violonĉelo} \textsuperscript{(83)}, \VUgras{trombono}
\textsuperscript{(84)}, \VUgras{tamburo} \textsuperscript{(85)}, k. c.. Oni
penas savi parton de la \VUgras{meblaro} \textsuperscript{(87)}.

%%K t11-c2-07-1 p
7. --- \VUgras{La aktoroj} \textsuperscript{(89)} kaj \VUgras{la aktorinoj}
\textsuperscript{(88)}, \VUgras{kantistoj} kaj \VUgras{kantistinoj} devis
forkuri kun sia teatra \VUgras{kostumo} \textsuperscript{(90)}. La tenoro
klarigas al \VUgras{ĵurnalisto} \textsuperscript{(92)} kiel tio okazis.
\VUgras{La raportisto} tuj alkuris de la unua momento kun la
aŭtoritatuloj : \VUgras{prefekto} \textsuperscript{(91)}, \VUgras{urbestro}
\textsuperscript{(93)}, \VUgras{komisaro de polico} \textsuperscript{(94)} kaj
\VUgras{juĝa oficisto} \textsuperscript{(95)}, kiuj enketos por juĝi pri la
respondecoj.
\VUsaut{6.0mm}
\VUfilet{22mm}
